\section{Related Work}
\label{sec:related}
\subsection{Coordinated Paths, Prediction, and Safety Control}
Cooperative planning must establish a compatible passing pattern and motions that vehicles can execute. Preiss et al.~\cite{preiss2017} combined discrete coordination with polynomial refinement and downwash-aware separation. H\"onig et al.~\cite{hoenig2018} connected collision-annotated roadmaps to continuous quadrotor trajectories, while CAPT~\cite{turpin2014} jointly selected assignments and trajectories for interchangeable robots. Discrete coordination can resolve the order of passage before continuous refinement. When assignments and arrival times are prescribed, the refinement must also respect the motion budget associated with each fixed endpoint. Changing an assignment or allocating more time is then outside the planning decision.

Distributed model predictive control (MPC) coordinates vehicles through predicted neighbor motions. Luis and Schoellig~\cite{luis2019} introduced on-demand collision constraints and penalized relaxation; Luis et al.~\cite{luis2020} developed online trajectory generation with continuous input parameterization and event-triggered replanning. Buffered Voronoi cells~\cite{zhou2017bvc} allocate local free space through relative-position geometry. Communication-efficient predictive formation control~\cite{schiz2025communication} reduces information exchange using autoencoders. These methods address local interaction and online update costs. Their rolling prediction windows differ from a prescribed completion horizon: the latter ties every remaining control interval to terminal position and velocity conditions that must be met at the same fixed time.

Control barrier functions (CBFs) constrain the current control action. Gon\c{c}alves et al.~\cite{goncalves2024} constructed safe drone acceleration fields with a quadratic program (QP), followed by thrust and angular-velocity tracking. Mestres et al.~\cite{mestre2024} distributed coupled safety constraints through auxiliary mismatch variables. These methods provide feasibility or safety results under their stated assumptions. Their control objectives concern instantaneous actions; a complete rest-to-rest motion with a fixed deadline requires additional temporal coordination. Their QP-based control updates provide a complementary level of coordination.

Deadlock studies further distinguish local safety from sustained progress. Grover et al.~\cite{grover2023} characterized reactive-controller deadlock through optimality conditions. Jankovic et al.~\cite{jankovic2024} linked gridlock propensity to equilibrium stability and developed predictor--corrector coordination. Chandra et al.~\cite{chandra2025} incorporated deadlock prevention into decentralized MPC with discrete-time CBFs. These results motivate explicit treatment of conflicting avoidance decisions. They also show why success under a fixed arrival time is a stronger task requirement than the availability of a safe instantaneous action or eventual progress toward a goal.

\subsection{Corridor-Based and Centralized Optimization}
Safe corridors turn geometric avoidance into conservative local constraints. Park and Kim~\cite{park2021} used relative safe flight corridors with relaxation and reciprocal avoidance for online quadrotor planning in dynamic environments. \c{S}enba\c{s}lar et al.~\cite{senbaslar2023} developed RLSS using linear spatial separations and polynomial QPs, with duration adjustment when derivative limits require it. Hard corridors are effective when their intersection with the motion constraints is nonempty. When that intersection is empty, a solver must be complemented by a change in local geometry, relaxation, or timing. Fixed-time tasks retain the first two options while excluding horizon extension.

Online planners address additional execution and communication requirements. EGO-Swarm~\cite{zhou2021ego} combines local optimization with perception and topological alternatives. MADER~\cite{tordesillas2022} uses asynchronous trajectory checking and commitment; Robust MADER~\cite{kondo2024} accounts for communication delay. PRIMER~\cite{kondo2025} learns perception-aware position--yaw planning while maintaining deconfliction. Park et al.~\cite{park2025} coordinate waypoints and subgoals with convergence and deadlock-prevention results under specified communication conditions. These capabilities are valuable for distributed operation. Centralized fixed-horizon planning instead couples the full joint motion in a known static map before execution.

Centralized formulations expose the full coupling among vehicles. Adaptive-scaling constrained optimization (ASCO)~\cite{li2021asco} changes the active collision-risk range during nonlinear cooperative planning. Priority-sharing intersection management~\cite{li2023priority} combines social-force references, optimization, and reservations. Sequential convex programming (SCP)~\cite{augugliaro2012} repeatedly updates local affine collision restrictions. Each approach uses a different response to difficult intermediate problems: active-set adjustment, coordination restructuring, or local approximation. In the fixed-team, fixed-endpoint problem considered here, intermediate failures must be handled while preserving the assigned task. A hard iterative corridor comparison is therefore especially informative: it tests whether reconstruction alone can provide the needed adjustment.

\subsection{From Iterative Corridors to Accepted Fixed-Time Motion}
Li et al.~\cite{li2022liom} introduced lightweight iterative optimization for parking, combining box corridors, softened nonlinear equalities, and repeated reconstruction. Bu et al.~\cite{bu2025} extended cooperative corridor planning with spatial and relative boxes, relative-corridor slack, kinematic penalties, and social-force initialization. These works establish the iterative corridor and penalty principles used here. The present formulation fixes each UAV's terminal position and zero terminal velocity at a common prescribed time. Its affine integration defects yield convex QPs, and reachable-speed padding connects nodal clearance to inter-node motion. A closed-form control-moment correction then enforces both terminal conditions before the continuous trajectory is checked.

The solution procedure gives a hard formulation a direct execution path whenever its returned motion passes the common checks. Unsuccessful entry or physical acceptance invokes relaxed reconstruction, tolerance adjustment, and exact finishing. These branches address distinct failures: unavailable local hard motions, residual kinematic mismatch, and changes introduced by output reconstruction. Matched hard iterative methods and independently checked conflict cases examine the role of each computational path. Table~\ref{tab:distinction} summarizes the methodological boundaries. For execution, linking planning to a tracking model is equally relevant; Huang et al.~\cite{huang2022srn} demonstrated this connection through smooth vehicle paths and control-oriented cosimulation. Here it is examined with a multirotor rigid-body simulation.

\begin{table}[htbp]
\begin{adjustwidth}{-\extralength}{0cm}
\caption{Representative formulations and their response to difficult local planning problems. The common completion deadline differs from a rolling prediction horizon.\label{tab:distinction}}
\small\setlength{\tabcolsep}{4pt}\renewcommand{\arraystretch}{1.13}
\begin{tabularx}{\fulllength}{p{2.8cm}YYY}
\toprule Method & Model and coordination & Timing / boundary commitment & Recovery or acceptance mechanism\\
\midrule
Li et al.~\cite{li2022liom} & Nonlinear parking model; spatial boxes & Optimized parking duration & Softened nonlinear equalities and iterative corridor reconstruction\\
Bu et al.~\cite{bu2025} & Cooperative vehicle model; STC/RSTC & Fixed traverse horizon; exit-region endpoints & Kinematic and relative-corridor penalties; repeated construction\\
Park and Kim~\cite{park2021} & Online multi-quadrotor relative corridors & Receding-horizon motion & Relative-corridor relaxation and reciprocal avoidance\\
RLSS~\cite{senbaslar2023} & Decentralized polynomial QPs & Adjustable segment durations & Linear separation and temporal rescaling for derivative limits\\
ASCO~\cite{li2021asco} & Centralized nonlinear cooperative planning & Optimized terminal time & Adaptation of the active collision-risk range\\
CBF control~\cite{goncalves2024,jankovic2024} & State-dependent safety-control QPs & Instantaneous control and navigation & Feasible safety filters and deadlock-aware policies under stated assumptions\\
SCP~\cite{augugliaro2012} & Centralized convex trajectory approximations & Prescribed trajectory duration & Updated affine separation constraints\\
\textbf{Proposed LIOM} & Affine multi-UAV QPs; spatial and relative regions & Fixed common time, positions, and rest endpoints & Hard entry, relaxed reconstruction, endpoint moment correction, and complete-interval acceptance\\
\bottomrule\end{tabularx}
\end{adjustwidth}
\end{table}
